All pins should now be enabled and allocated. In the GPIO section, rename the pins to match the following labels:
\begin{table}[H]
    \centering
    \begin{tabular}{ll}
        \hline
        \textbf{Pin} & \textbf{Name} \\
        \hline
        PB3 & RESETBAR \\
        PB4 & TRIGGER \\
        PB6 & SHUTTER \\
        PB14 & MCU\_LED\_1 \\
        PB15 & MCU\_LED\_2 \\
        PC0 & PWR\_STATE \\
        PC1 & PWR\_ACTION \\
        PC13 & SCROLL\_B \\
        PD8 & READ\_LED\_1 \\
        PD9 & READ\_LED\_2 \\
        PD10 & WRITE\_LED\_1 \\
        PD11 & WRITE\_LED\_2 \\
        PD12 & DEBUG\_LED \\
        PE2 & SCROLL\_A \\
        PE13 & SENSOR\_LED\_2 \\
        PE14 & SENSOR\_LED\_1 \\
        PG15 & STANDBY \\
        PH3-BOOT0 & BOOT0 \\
        PH9 & PIN81 \\
        PH10 & PIN84 \\
        PH11 & PIN85 \\
        PI3 & PIN97 \\
        \hline
    \end{tabular}
    \caption{List of user pin labels for C2Sv0.1 STM32CubeIDE Project.}
\end{table}

Now we must ensure that the \verb|SCROLL_A| and \verb|SCROLL_B| pins (PC13 and PE2) are pulled up for the rotary encoder code to function.